\section*{الفصل \ref{ch:b1:stratigraphic-scale} --- السلّم الطبقي}

\begin{solution}{exo:b1:stratigraphic-scale:1}
(أ)~كرونوستراتغرافي، نظام. (ب)~صخري، تكوين.
(ج)~جيوكرونولوجي، حقبة (وصخرها هو \omterm{def:b1:stratigraphic-scale:chronostratigraphy}{السلسلة} الكرتاسية السفلى).
(د)~بيوستراتغرافي، حيّز أحيائي. (ه)~كرونوستراتغرافي، طابق.
(و)~جيوكرونولوجي، حقبة. (ز)~صخري، عضو.
\end{solution}

\begin{solution}{exo:b1:stratigraphic-scale:2}
الباليوزويك $538.8 - 251.902 = \qty{286.9}{\Myr}$، أي \qty{53.2}{\%} من الفانيروزويك؛ الميزوزويك $251.902 - 66.00 = \qty{185.9}{\Myr}$، أي \qty{34.5}{\%}؛ السينوزويك \qty{66.0}{\Myr}، أي \qty{12.2}{\%}. والفانيروزويك هو $538.8/4567 = \qty{11.8}{\%}$ من الزمن منذ أساس الهاديان.
\end{solution}

\begin{solution}{exo:b1:stratigraphic-scale:3}
تكوّنت A وB وC وD بهذا الترتيب (التراكب)، أفقيةً (الأفقية الأصلية)، فميلها حدث لاحق. والدايك E يقطع A إلى D، فهو أحدث من D (علاقات القطع)؛ ولا يمكن تحديد إن كان جاء قبل الميل أو بعده من الوصف. وسطح التآكل يقطع D وE، فهو أحدث منهما (علاقات القطع). ثم رُسّبت عليه F وG (التراكب). والصدع يزيح A إلى F دون G: فقد تحرّك بعد F وقبل G (علاقات القطع). الترتيب: A، B، C، D؛ والميل ودخول E (الترتيب غير محسوم)؛ والتآكل؛ وF؛ والتصدّع؛ وG.
\end{solution}

\begin{solution}{exo:b1:stratigraphic-scale:4}
$r = 30/(85.40 - 84.60) = \qty{37.5}{m}$ لكل مليون سنة. ويقع الظهور الأول على بُعد $12/37.5 = \qty{0.32}{\Myr}$ بعد الرماد الأسفل: $85.40 - 0.32 =
\qty{85.08}{\Ma}$.
\end{solution}

\begin{solution}{exo:b1:stratigraphic-scale:5}
(أ)~أربعة أعمدة رأسية مقابل الارتفاع. (ب)~يظهر B وC معًا من \qtyrange{26}{50}{m}: حيّز مدى متزامن سماكته \qty{24}{m}. \omterm{def:b1:stratigraphic-scale:biozone}{والحيّز الفاصل} من الظهور الأول لB إلى الظهور الأول لD يمتدّ من \qtyrange{12}{41}{m}، أي \qty{29}{m}. (ج)~في القطاع الثاني يمكن تحديد الظهورين الأولين لB (\qty{3}{m}) ولC (\qty{9}{m})، فيُعرف أساس \omterm{def:b1:stratigraphic-scale:biozone}{الحيّز الفاصل} في (ب) وأساس \omterm{def:b1:stratigraphic-scale:biozone}{حيّز المدى} المتزامن؛ أما قمّته (الظهور الأول لD) وقمّة التداخل كما تُرى في القطاع~1 فلا يمكن فحصهما لغياب D. والظهوران الأوّلان لB وC ($\qty{12}{m} \leftrightarrow
\qty{3}{m}$، $\qty{26}{m} \leftrightarrow \qty{9}{m}$) هما أكثر الخطوط أمانًا؛ أما الظهورات الأخيرة (لB عند \qty{50}{m} و\qty{15}{m}) فتعتمد أكثر على الحفظ.
\end{solution}

\begin{solution}{exo:b1:stratigraphic-scale:6}
(أ)~$60/20 = \qty{3}{\Myr}$: قمّة الحجر الرملي أحدث بثلاثة ملايين سنة في البئر البرّي. (ب)~مليون سنة من الطين بمعدّل \qty{30}{m} لكل مليون سنة: يقع الرماد على \qty{30}{m} فوق قمّة الحجر الرملي. (ج)~حين سقط الرماد، كان الشاطئ قد تحرّك \qty{20}{km} نحو اليابسة من البئر البحري، فكان ما زال على بُعد \qty{40}{km} من البحر عن البئر البرّي، الذي كان واقفًا على اليابسة. سقط الرماد هناك على سطح اليابسة، تحت الحجر الرملي الذي رسّبه البحر بعد ذلك بمليوني سنة، إن لم يُغسل قبل ذلك، كما تُغسل طبقة مفكّكة على اليابسة غالبًا. وطبقة الرماد نفسها فوق حدّ \omterm{def:b1:stratigraphic-scale:lithostratigraphy}{التكوين} في بئر، وتحته في الآخر: فالحدّ يعبر الزمن.
\end{solution}

\begin{solution}{exo:b1:stratigraphic-scale:7}
(أ)~اللوكوفيان $1.36/419.62 = \qty{0.32}{\%}$؛ الإنديان $0.024/251.902 =
\qty{0.0095}{\%}$، أي أصغر بنحو 34 مرّة. وقاعدة الإنديان تقع بين طبقات رماد مؤرّخة بدقّة في القطاع نفسه للGSSP؛ وعدم يقين كبير يوحي بحدّ يجب أن يُحمل عمره من مستويات مؤرّخة أبعد. (ب)~$251.941 - 251.880 = \qty{0.061}{\Myr} = 61$ ألف سنة، مع عدم يقين $\sqrt{37^2 + 31^2} = 48$ ألف سنة: \num[uncertainty-mode=separate]{60 \pm 48} ألف سنة المنشورة.
\end{solution}

\begin{solution}{exo:b1:stratigraphic-scale:8}
(أ)~$m = 3000 \times \frac43\pi\,(5000)^3 = \qty{1.57e15}{kg}$؛ والغبار، $60m = \qty{9.42e16}{kg}$. (ب)~مساحة سطح الأرض $4\pi(6.378\times10^6)^2
= \qty{5.11e14}{m^2}$: \qty{184}{kg} لكل متر مربّع. (ج)~$184/2500 =
\qty{0.074}{m}$، نحو \qty{7}{cm}. (د)~تتناسب السماكة مع مكعب القطر: \qty{1.6}{cm} ل\qty{6}{km}، و\qty{20}{cm} ل\qty{14}{km}. وترتيب المقدار، من السنتيمترات إلى الديسيمترات، هو ترتيب طين الحدّ البالغ سمكه \qty{50}{cm} في الكاف؛ والمقارنة لا يمكن أن تكون إلا تقريبية، إذ لا يمكن أن يسقط الغبار بالتساوي، والطين يجمع أيضًا الرسوب المعتاد في زمنه.
\end{solution}

\begin{solution}{exo:b1:stratigraphic-scale:9}
(أ)~$\num{60e3} \times \num{2e-6} = \qty{0.12}{m}$، أي \qty{12}{cm}؛ ول\num{12} و\num{108} ألف سنة، \qty{2.4}{cm} و\qty{21.6}{cm}. (ب)~$\num{32e3} \times
\num{10e-6} = \qty{0.32}{m}$. (ج)~يشغل الانقراض سنتيمترات إلى ديسيمترات من الصخر: فيجب أخذ العيّنات سنتيمترًا بسنتيمتر، وعيّنة واحدة سماكتها \qty{10}{cm} في القطاع المكثّف تتوسّط خمسين ألف سنة، أي معظم الحدث.
\end{solution}

\begin{solution}{exo:b1:stratigraphic-scale:10}
(أ)~$u = R/s = 10 \times 2000 = \qty{20000}{m}$ لكل مليون سنة، أي \qty{20}{km} لكل مليون سنة؛ وعلى \qty{30}{km}، $\Delta t = \qty{1.5}{\Myr}$. (ب)~أصبح السهل الآن السطح $z = s x + a t$. ويحقّق الشاطئ $s x + a t = z_0 + R t$، ف$x(t) = (z_0 + (R - a)t)/s$ و$\mathrm dx/\mathrm dt = (R - a)/s$، موجبة (نحو اليابسة) ما دام $a < R$. وحين يكون $a > R$ تكون السرعة سالبة: ينتقل الشاطئ نحو البحر، وهو تراجع بحري، ويصير حجر الشاطئ الرملي حينئذٍ أحدث نحو البحر. (ج)~جرف طوله \qty{200}{m} يمتدّ، بسرعة \qty{20}{km} لكل مليون سنة، على $0.2/20 = \qty{0.01}{\Myr}$: عشرة آلاف سنة من التباين الزمني، أقلّ بكثير مما تستطيع الحفريات أو معظم التأريخات تمييزه؛ والصخر على جانبي الحدّ هو نفسه في كل مكان على امتداده.
\end{solution}

\begin{solution}{exo:b1:stratigraphic-scale:11}
(أ)~$f = (25 - 10)/(50 - 10) = 0.375$ و$t = 100.50 - 0.375 \times 0.80 =
\qty{100.20}{\Ma}$، مع
\[
  \sigma_t = \sqrt{(0.625 \times 0.05)^2 + (0.375 \times 0.03)^2} = \qty{0.033}{\Myr}.
\] (ب)~$g(f) = (1 - f)^2\sigma_1^2 + f^2\sigma_2^2$
له $g'(f) = -2(1 - f)\sigma_1^2 + 2f\sigma_2^2$، وصفره عند $f^* =
\sigma_1^2/(\sigma_1^2 + \sigma_2^2) = 0.0025/0.0034 = 0.735$، أقرب إلى الطبقة المؤرّخة الأفضل؛ وعندها $\sigma_t = \sigma_1\sigma_2/\sqrt{\sigma_1^2 + \sigma_2^2} =
\qty{0.026}{\Myr}$، أصغر من كليهما. (ج)~$g$ قطع مكافئ مفتوح نحو الأعلى، فعلى $[0, 1]$ يبلغ أكبر قيمة عند أحد الطرفين: $g(0) = \sigma_1^2$، و$g(1) = \sigma_2^2$. ومن ثمّ $\sigma_t \le \max(\sigma_1, \sigma_2)$ بين الطبقتين.
\end{solution}

\begin{solution}{exo:b1:stratigraphic-scale:12}
(أ)~الحدّ في منتصف العمق: $f = 0.5$، \qty{29.50}{\Ma}. (ب)~مع $\varphi_0\lambda = 0.63 \times 1960 = \qty{1234.8}{m}$، فالسماكة الصلبة من \qtyrange{600}{1000}{m} هي $400 - 1234.8\,(e^{-600/1960} - e^{-1000/1960}) =
\qty{232.2}{m}$، ومن \qtyrange{200}{600}{m}، \qty{194.2}{m}؛ والمجموع \qty{426.3}{m}. وكسر الحبيبات المرسّبة بين الرماد الأسفل والحدّ هو $232.2/426.3 = 0.545$، والعمر $30.0 - 0.545 =
\qty{29.46}{\Ma}$. (ج)~التقدير الثاني أصغر ب\num{45} ألف سنة، وهو الأفضل: فالنصف الأعمق فقد ماء أكثر، فيحمل حبيبات، وزمنًا، أكثر لكل متر من النصف الأضحل (مسامية متوسطة \num{0.42} مقابل \num{0.51}).
\end{solution}

\subsection*{المسألة: ترابط ثلاثة قطاعات عبر حوض}

\begin{solution}{pb:b1:stratigraphic-scale:1}
\begin{enumerate}
  \item ثلاثة أعمدة بمقياس \qty{1}{cm} لكل \qty{10}{m} تقريبًا، موصولة بطبقات الرماد وبحدود الحيّزات؛ وإسفين حجر كِل الرملي يتناقص من \qty{45}{m} في A إلى \qty{6}{m} في C.
  \item الحيّز~2: في A، $65.8 - 36.6 = \qty{29.2}{m}$، وفي B \qty{20.8}{m}، وفي C \qty{10.0}{m}. والحيّز~3 المسجَّل: في A \qty{14.2}{m}، وفي B \qty{14.0}{m}، وفي C \qty{7.2}{m}.
  \item يبدأ الحيّز~1 عند الظهور الأول ل\emph{V.~primum}، الذي يوجد من قاعدة كل قطاع: فظهوره يقع تحت القطاعات كلها.
  \item إنه يعيش في الحجر الرملي وحده: فهو يتبع السحنة، والسحنة أعمارها تختلف من مكان إلى آخر. وفي A يوجد في الحيّزين~1 و\num{2}، وفي B وC في الحيّز~1 وحده.
  \item في الحيّز~2 في الثلاثة كلها (A \qty{49.0}{m}، وB \qty{40.0}{m}، وC \qty{19.0}{m}). وهو متّسق من قطاع إلى آخر هنا، فيمكن أن يقسّم الحيّز~2، لكن الظهور الأخير يعتمد على الحفظ، وينبغي فحصه في قطاعات أكثر أولًا.
  \item تصل الخطوط \qty{36.6}{m} و\qty{31.2}{m} و\qty{14.8}{m}، ثم \qty{65.8}{m} و\qty{52.0}{m} و\qty{24.8}{m}. ويرقّ الحيّز~2 بنسبة $29.2 : 20.8 : 10.0$، أي نحو $2.9 : 2.1 : 1$: فلو دام الحيّز الزمن نفسه في كل مكان، لتراكم الرسوب في A نحو ثلاث مرّات أسرع منه في C.
  \item A: الحيّز~2 ($45.0 > 36.6$)؛ B: الحيّز~1 ($20.0 < 31.2$)؛ C: الحيّز~1 ($6.0 < 14.8$).
  \item لا: فهو في الحيّز~2 في A، وفي الحيّز~1 في B وC، فهو أحدث نحو الغرب. وانتشر الطين البحري نحو الغرب فوق رمل الشاطئ: تحرّك الشاطئ نحو الغرب، أي نحو اليابسة، في تجاوز بحري.
  \item يسقط الرماد على الحوض كله في لحظة جيولوجية، ويستقرّ على أيّ رسوب؛ وقمّة الحجر الرملي تتبع الشاطئ المتحرّك، وتكون متباينة الزمن.
  \item استمرار الطبقة جانبيًّا (الطبقة تمتدّ عبر الحوض)، والتراكب (الرماد~1 تحت الرماد~2 في كل قطاع، وهذا يتحقّق من الترابط).
  \item A: $(71.0 - 5.0)/1.20 = \qty{55.0}{m}$ لكل مليون سنة؛ B: $48.0/1.20 = \qty{40.0}{}$؛ C: $24.0/1.20 = \qty{20.0}{}$.
  \item A هي الأقرب إلى اليابسة، حيث أُمدّ الرسوب؛ ووصل إلى الماء العميق نحو الشرق منه أقلّ.
  \item $0.1/55 = \num{1.8}$ ألف سنة في A، و$\num{2.5}$ في B، و$\num{5.0}$ في C: فA تُميّز الزمن بأدقّ قدر.
  \item A: $152.40 - 40.0/55.0 = \qty{151.67}{\Ma}$؛ B: $152.40 - 12.0/40.0 =
        \qty{152.10}{\Ma}$؛ C: $152.40 - 3.0/20.0 = \qty{152.25}{\Ma}$.
  \item من C إلى B: $20/(152.25 - 152.10) = \qty{133}{km}$ لكل مليون سنة؛ ومن B إلى A: $25/(152.10 - 151.67) = \qty{58.5}{km}$ لكل مليون سنة. فقد تباطأ التجاوز.
  \item A: $152.40 - 60.8/55.0 = \qty{151.29}{\Ma}$؛ B: $152.40 - 44.0/40.0 =
        \qty{151.30}{\Ma}$؛ C: $152.40 - 21.8/20.0 = \qty{151.31}{\Ma}$.
  \item تمتدّ على \num{16} ألف سنة، وهذا أقلّ من عدم يقين كل منها (نحو \num{28} ألف سنة، السؤال~19): فالظهور الأول متزامن عبر الحوض ضمن دقّة التأريخات، وهذا ما يجعله دليلًا جيدًا.
  \item A: $152.40 - 31.6/55.0 = \qty{151.826}{\Ma}$؛ B: $152.40 - 23.2/40.0 =
        \qty{151.820}{\Ma}$؛ C: $152.40 - 11.8/20.0 = \qty{151.810}{\Ma}$.
  \item $f = 0.479$، و$0.483$، و$0.492$؛ و$\sigma_t = \sqrt{(1 - f)^2 + f^2} \times
        0.04 = \qty{0.028}{\Myr}$ في كل منها، أي \num{28} ألف سنة.
  \item المتوسط \qty{151.819}{\Ma}؛ والانتشار $151.826 - 151.810 = \num{16}$ ألف سنة، وهو أصغر من عدم يقين كل عمر: فالثلاثة تتفق.
  \item يفترض الاستكمال معدّلًا واحدًا من الرماد~1 إلى الرماد~2. وفي A يحوي ذلك الفاصل حجرًا رمليًّا شاطئيًّا حتى \qty{45}{m} وصخرًا طينيًّا فوقه؛ ولا يلزم أن يكون الرمل والطين قد تراكما بالمعدّل نفسه، ولا أن ينضغطا بالطريقة نفسها (\cref{ch:b1:diagenesis-basins}). فإن تراكم الرمل أسرع، فجزء الحجر الرملي يحمل زمنًا أقلّ مما توحي به سماكته، ويكون الظهور الأول أقدم مما حُسب، والعمر من A هو فعلًا الأقدم بين الثلاثة.
  \item B: يقع الظهور الأول في صخر طيني بحري، على \qty{11}{m} فوق قمّة الحجر الرملي، دون تغيّر في السحنة قرب المستوى، ويعطي المعدّل \num{2.5} ألف سنة لكل \qty{10}{cm}. وA تُميّز الزمن بشكل أفضل، لكنها تضع المستوى في الحجر الرملي الشاطئي، على بُعد \qty{8.4}{m} تحت تغيّر في السحنة؛ وC مكثّفة، \num{5} ألف سنة لكل \qty{10}{cm}. فC هي الأسوأ في التمييز، وA الأسوأ في السحنة.
  \item من متوسطي عمري الظهورين الأولين: $151.819 - 151.302 =
        \qty{0.52}{\Myr}$: نحو نصف مليون سنة.
  \item الحدّ هو \qty[uncertainty-mode=separate]{151.82(3)}{\Ma}.
\end{enumerate}
\end{solution}
